\documentclass[10pt,a4paper]{amsart}
\usepackage{amsmath,amssymb,amsthm,graphicx}
% Retained preamble for the Stacks Project chapter divisors.tex (Divisors).
% The chapter file pinned-original-source.tex carries only \input{preamble} on its first line, so no
% preamble text lives inside the chapter itself.  The Stacks Project uses one single shared preamble for
% every chapter; the snapshot stages it verbatim in sources/stacks-chow/preamble.tex.  The block below is
% copied mechanically, line for line, from that file's lines 14-18 (the Xy-pic setup) and lines 154-217
% (the native theorem declarations together with the mathematical macros).  Nothing was added, removed or
% redefined.  The wrapper class and the amsmath/amssymb/amsthm/graphicx base packages are supplied by the
% assembler, so the class/\externaldocument part of the shared preamble is not repeated here.
% For commutative diagrams we use Xy-pic
\usepackage[all]{xy}

% We use 2cell for 2-commutative diagrams.
\xyoption{2cell}
% Theorem environments.
%
\theoremstyle{plain}
\newtheorem{theorem}[subsection]{Theorem}
\newtheorem{proposition}[subsection]{Proposition}
\newtheorem{lemma}[subsection]{Lemma}

\theoremstyle{definition}
\newtheorem{definition}[subsection]{Definition}
\newtheorem{example}[subsection]{Example}
\newtheorem{exercise}[subsection]{Exercise}
\newtheorem{situation}[subsection]{Situation}

\theoremstyle{remark}
\newtheorem{remark}[subsection]{Remark}
\newtheorem{remarks}[subsection]{Remarks}

\numberwithin{equation}{subsection}

% Macros
%
\def\lim{\mathop{\mathrm{lim}}\nolimits}
\def\colim{\mathop{\mathrm{colim}}\nolimits}
\def\Spec{\mathop{\mathrm{Spec}}}
\def\Hom{\mathop{\mathrm{Hom}}\nolimits}
\def\Ext{\mathop{\mathrm{Ext}}\nolimits}
\def\SheafHom{\mathop{\mathcal{H}\!\mathit{om}}\nolimits}
\def\SheafExt{\mathop{\mathcal{E}\!\mathit{xt}}\nolimits}
\def\Sch{\mathit{Sch}}
\def\Mor{\mathop{\mathrm{Mor}}\nolimits}
\def\Ob{\mathop{\mathrm{Ob}}\nolimits}
\def\Sh{\mathop{\mathit{Sh}}\nolimits}
\def\NL{\mathop{N\!L}\nolimits}
\def\CH{\mathop{\mathrm{CH}}\nolimits}
\def\proetale{{pro\text{-}\acute{e}tale}}
\def\etale{{\acute{e}tale}}
\def\QCoh{\mathit{QCoh}}
\def\Ker{\mathop{\mathrm{Ker}}}
\def\Im{\mathop{\mathrm{Im}}}
\def\Coker{\mathop{\mathrm{Coker}}}
\def\Coim{\mathop{\mathrm{Coim}}}

% Boxtimes
%
\DeclareMathSymbol{\boxtimes}{\mathbin}{AMSa}{"02}

%
% Macros for moduli stacks/spaces
%
\def\QCohstack{\mathcal{QC}\!\mathit{oh}}
\def\Cohstack{\mathcal{C}\!\mathit{oh}}
\def\Spacesstack{\mathcal{S}\!\mathit{paces}}
\def\Quotfunctor{\mathrm{Quot}}
\def\Hilbfunctor{\mathrm{Hilb}}
\def\Curvesstack{\mathcal{C}\!\mathit{urves}}
\def\Polarizedstack{\mathcal{P}\!\mathit{olarized}}
\def\Complexesstack{\mathcal{C}\!\mathit{omplexes}}
% \Pic is the operator that assigns to X its picard group, usage \Pic(X)
% \Picardstack_{X/B} denotes the Picard stack of X over B
% \Picardfunctor_{X/B} denotes the Picard functor of X over B
\def\Pic{\mathop{\mathrm{Pic}}\nolimits}
\def\Picardstack{\mathcal{P}\!\mathit{ic}}
\def\Picardfunctor{\mathrm{Pic}}
\def\Deformationcategory{\mathcal{D}\!\mathit{ef}}

% --- appended 2026-10-12 (Chief): the project own definitions of the reference/slogan/history
% display environments, copied mechanically from the shared Stacks preamble lines 8-12
% (\usepackage{verbatim} plus the three \newenvironment lines, which the chapter preamble keeps
% verbatim from the project preamble). Nothing added, removed or redefined: these environments
% carry no mathematical content and the project itself typesets them as comments. Without these
% three lines a counted statement whose own text uses them fails with Environment slogan undefined.
\usepackage{verbatim}
\newenvironment{reference}{\comment}{\endcomment}
\newenvironment{slogan}{\comment}{\endcomment}
\newenvironment{history}{\comment}{\endcomment}
\usepackage{cleveref}
\crefname{originaldefinition}{original definition}{original definitions}
\Crefname{originaldefinition}{Original definition}{Original definitions}
\makeatletter
\newlabel{definition-invertible-sheaf-effective-Cartier-divisor}{{01WX}{source}{}{sourcepointer-1}{}}
\newlabel{definition-invertible-sheaf-effective-Cartier-divisor@cref}{{[originaldefinition][0][]01WX}{[1][0][]source}{}{}{}}
\makeatother
\crefname{originallemma}{original lemma}{original lemmas}
\Crefname{originallemma}{Original lemma}{Original lemmas}
\makeatletter
\newlabel{lemma-complement-open-affine-effective-cartier-divisor-bis}{{0EGJ}{source}{}{sourcepointer-2}{}}
\newlabel{lemma-complement-open-affine-effective-cartier-divisor-bis@cref}{{[originallemma][0][]0EGJ}{[1][0][]source}{}{}{}}
\makeatother
\crefname{originallemma}{original lemma}{original lemmas}
\Crefname{originallemma}{Original lemma}{Original lemmas}
\makeatletter
\newlabel{more-algebra-lemma-regular-local-UFD}{{0AG0}{source}{}{sourcepointer-3}{}}
\newlabel{more-algebra-lemma-regular-local-UFD@cref}{{[originallemma][0][]0AG0}{[1][0][]source}{}{}{}}
\makeatother
\crefname{originaldefinition}{original definition}{original definitions}
\Crefname{originaldefinition}{Original definition}{Original definitions}
\makeatletter
\newlabel{morphisms-definition-family-ample-invertible-modules}{{0FXR}{source}{}{sourcepointer-4}{}}
\newlabel{morphisms-definition-family-ample-invertible-modules@cref}{{[originaldefinition][0][]0FXR}{[1][0][]source}{}{}{}}
\makeatother
\crefname{originallemma}{original lemma}{original lemmas}
\Crefname{originallemma}{Original lemma}{Original lemmas}
\makeatletter
\newlabel{properties-lemma-normal-Noetherian}{{033M}{source}{}{sourcepointer-5}{}}
\newlabel{properties-lemma-normal-Noetherian@cref}{{[originallemma][0][]033M}{[1][0][]source}{}{}{}}
\makeatother
\crefname{originallemma}{original lemma}{original lemmas}
\Crefname{originallemma}{Original lemma}{Original lemmas}
\makeatletter
\newlabel{properties-lemma-regular-normal}{{0567}{source}{}{sourcepointer-6}{}}
\newlabel{properties-lemma-regular-normal@cref}{{[originallemma][0][]0567}{[1][0][]source}{}{}{}}
\makeatother
\title{A quasi-compact regular scheme with affine diagonal has an ample family of invertible modules}
\author{The Stacks Project Authors}
\date{stacks-project:divisors.tex:local-snapshot-20261012; GNU FDL 1.2 or later}
\begin{document}
\maketitle
\noindent\textit{Editorial note.} All delivered mathematical text is the author's own, copied from the pinned Stacks Project chapter file pinned-original-source.tex; no statement or proof step was rewritten, reordered or removed.\par\medskip
\noindent\textit{Editorial note.} Following the multi-claim convention for this source, source\_url is the source-level URL https://stacks.math.columbia.edu/ and the source label lemma-regular-ample-family is carried in claim\_key and here rather than in a per-claim tag URL.\par\medskip
\noindent\textit{Editorial note.} The preamble is the chapter's shared Stacks Project preamble, retained as sources/stacks-divisors/retained-preamble.tex: a line-for-line excerpt of the shared Stacks preamble staged with this snapshot. No macro was invented or redefined, and this card adds no preamble text of its own.\par\medskip
\noindent\textit{Editorial note.} Statement span 3553-3559 and proof span 3561-3581 come from the halfA pool card operations/pool-cards/stacks-divisors-halfA.json and were re-read in the pinned chapter file before this card was written.\par\medskip
\noindent\textit{Editorial note.} The statement refers to the definition of an ample family of invertible modules in another chapter; the reference is part of the counted statement span and is handled as a navigation pointer.\par\medskip
\noindent\textit{Editorial note.} 6 cross-references occur inside the counted spans; they are listed in external\_references as navigation pointers only.\par\medskip
\noindent\textit{Allowed standard prerequisites (editorial).}\begin{itemize}
\item A regular scheme is a finite disjoint union of integral schemes, and a normal Noetherian scheme has finitely many irreducible components, cited from the properties chapter. Cited at proof lines 3561-3564. Reduces the statement to the integral case.
\item Local rings of a regular scheme are unique factorisation domains, cited from the more-algebra chapter. Cited at proof lines 3569-3570. Makes the complement-of-an-affine-open statement applicable.
\item For a quasi-compact scheme with affine diagonal and UFD local rings, the complement of an affine open is the support of an effective Cartier divisor, cited from the same chapter. Cited at proof lines 3570-3572. Produces the divisor whose associated invertible module is the candidate of the ample family.
\item The invertible module associated with an effective Cartier divisor and its canonical section, whose nonvanishing locus is the complement of the divisor, cited from the same chapter. Cited at proof lines 3573-3577. Identifies the nonvanishing locus of the canonical section with the chosen affine open, which is the cap condition for an ample family.
\end{itemize}\medskip
\ifdefined\hypertarget\hypertarget{sourcepointer-1}{}\fi
\noindent\textit{Source reference.} Original definition 01WX: focused-seeds/AGP/sources/stacks-snapshot-20261011/tags.txt, source lines 2498--2498. Same-chapter reference used in the retained author text at source line(s) 3576, 6928, 9257. It is supplied only as a navigation pointer so that the retained text compiles; the counted item delivers the author's own statement and the author's complete proof, and no part of the referenced item is copied in or substituted for the counted argument.\par\medskip
\ifdefined\hypertarget\hypertarget{sourcepointer-2}{}\fi
\noindent\textit{Source reference.} Original lemma 0EGJ: focused-seeds/AGP/sources/stacks-snapshot-20261011/tags.txt, source lines 18099--18099. Same-chapter reference used in the retained author text at source line(s) 3572. It is supplied only as a navigation pointer so that the retained text compiles; the counted item delivers the author's own statement and the author's complete proof, and no part of the referenced item is copied in or substituted for the counted argument.\par\medskip
\ifdefined\hypertarget\hypertarget{sourcepointer-3}{}\fi
\noindent\textit{Source reference.} Original lemma 0AG0: focused-seeds/AGP/sources/stacks-snapshot-20261011/tags.txt, source lines 13178--13178. Same-chapter reference used in the retained author text at source line(s) 3149, 3570. It is supplied only as a navigation pointer so that the retained text compiles; the counted item delivers the author's own statement and the author's complete proof, and no part of the referenced item is copied in or substituted for the counted argument.\par\medskip
\ifdefined\hypertarget\hypertarget{sourcepointer-4}{}\fi
\noindent\textit{Source reference.} Original definition 0FXR: focused-seeds/AGP/sources/stacks-snapshot-20261011/tags.txt, source lines 19891--19891. Same-chapter reference used in the retained author text at source line(s) 3558, 3579. It is supplied only as a navigation pointer so that the retained text compiles; the counted item delivers the author's own statement and the author's complete proof, and no part of the referenced item is copied in or substituted for the counted argument.\par\medskip
\ifdefined\hypertarget\hypertarget{sourcepointer-5}{}\fi
\noindent\textit{Source reference.} Original lemma 033M: focused-seeds/AGP/sources/stacks-snapshot-20261011/tags.txt, source lines 4061--4061. Same-chapter reference used in the retained author text at source line(s) 3565. It is supplied only as a navigation pointer so that the retained text compiles; the counted item delivers the author's own statement and the author's complete proof, and no part of the referenced item is copied in or substituted for the counted argument.\par\medskip
\ifdefined\hypertarget\hypertarget{sourcepointer-6}{}\fi
\noindent\textit{Source reference.} Original lemma 0567: focused-seeds/AGP/sources/stacks-snapshot-20261011/tags.txt, source lines 6700--6700. Same-chapter reference used in the retained author text at source line(s) 3564. It is supplied only as a navigation pointer so that the retained text compiles; the counted item delivers the author's own statement and the author's complete proof, and no part of the referenced item is copied in or substituted for the counted argument.\par\medskip
\begin{lemma}
\label{lemma-regular-ample-family}
Let $X$ be a quasi-compact, regular scheme with affine diagonal.
Then $X$ has an ample family of invertible modules
(Morphisms, Definition
\ref{morphisms-definition-family-ample-invertible-modules}).
\end{lemma}

\begin{proof}
Observe that $X$ is a finite disjoint union of integral schemes
(Properties, Lemmas
\ref{properties-lemma-regular-normal} and
\ref{properties-lemma-normal-Noetherian}).
Thus we may assume that $X$ is integral as well as Noetherian,
regular, and having affine diagonal. Let $x \in X$.
Choose an affine open neighbourhood $U \subset X$ of $x$.
Since $X$ is integral, $U$ is dense in $X$.
By More on Algebra, Lemma \ref{more-algebra-lemma-regular-local-UFD}
the local rings of $X$ are UFDs. Hence by Lemma
\ref{lemma-complement-open-affine-effective-cartier-divisor-bis}
we can find an effective Cartier divisor $D \subset X$
whose complement is $U$. Then the canonical section
$s = 1_D$ of $\mathcal{L} = \mathcal{O}_X(D)$,
see Definition \ref{definition-invertible-sheaf-effective-Cartier-divisor},
vanishes exactly along $D$ hence $U = X_s$.
Thus both conditions in Morphisms, Definition
\ref{morphisms-definition-family-ample-invertible-modules}
hold and we are done.
\end{proof}

\end{document}
